\documentclass[11pt,a4paper]{article}

\usepackage[utf8]{inputenc}
\usepackage[T1]{fontenc}
\usepackage{lmodern}

\usepackage{geometry}
\geometry{margin=2.5cm}

\usepackage{hyperref}
\usepackage{xcolor}
\usepackage{graphicx}
\usepackage{listings}
\usepackage{booktabs}
\usepackage{longtable}
\usepackage{amsmath}
\usepackage{amssymb}

\hypersetup{
    colorlinks=true,
    linkcolor=blue,
    urlcolor=blue
}

\title{
\textbf{ClimateProcessingEngine}\\
\large A Professional Python Framework for Climate Data Processing and Statistical Analysis
}

\author{
Amin Fazl Kazemi\\
\href{https://github.com/AminFazlKazemi}{GitHub}
}

\date{\today}


\begin{document}

\maketitle


\begin{center}

\href{https://github.com/AminFazlKazemi/ClimateProcessingEngine}{
\textcolor{blue}{
github.com/AminFazlKazemi/ClimateProcessingEngine
}
}

\end{center}


\section*{Overview}


\textbf{ClimateProcessingEngine} is an open-source Python framework designed for
large-scale climate data processing, statistical distribution fitting,
geospatial analysis, and scientific visualization.


The framework provides an integrated workflow for:

\begin{itemize}

\item Climate time-series processing
\item Temperature and precipitation analysis
\item Missing data handling
\item Statistical distribution fitting
\item Climate variability assessment
\item Spatial analysis
\item Scientific visualization
\item Large-scale Zarr data processing

\end{itemize}


The project is developed using the scientific Python ecosystem:

\begin{itemize}

\item NumPy
\item Pandas
\item Xarray
\item Dask
\item SciPy
\item Scikit-learn
\item Rasterio
\item GeoPandas
\item Cartopy
\item Matplotlib
\item Numba

\end{itemize}



\section{Target Audience}


ClimateProcessingEngine is designed for:

\begin{itemize}

\item Climate scientists
\item Meteorologists
\item Hydrologists
\item Environmental researchers
\item Agricultural scientists
\item Geospatial data engineers
\item Graduate students

\end{itemize}



\section{Key Features}


\subsection{Climate Data Processing}


The engine supports:


\begin{itemize}

\item Minimum temperature analysis ($t_{min}$)
\item Mean temperature analysis ($t_{mean}$)
\item Maximum temperature analysis ($t_{max}$)
\item Precipitation processing
\item Climate index calculation
\item Long-term variability analysis

\end{itemize}



\subsection{Quality Control}


Implemented capabilities:


\begin{itemize}

\item Missing value detection
\item Outlier identification
\item Data validation
\item Statistical preprocessing

\end{itemize}



\subsection{Missing Data Handling}


Advanced approaches include:


\begin{itemize}

\item Window-based estimation
\item Interpolation techniques
\item Robust statistical methods

\end{itemize}



\section{Geospatial Analysis}


Supported operations:


\begin{itemize}

\item Raster processing
\item Vector processing
\item Coordinate transformation
\item Spatial interpolation
\item Regional climate mapping
\item Scientific cartography

\end{itemize}


Core libraries:


\begin{itemize}

\item Rasterio
\item GeoPandas
\item Xarray
\item Cartopy
\item Shapely

\end{itemize}
\section{Architecture}


ClimateProcessingEngine follows a modular pipeline-based architecture.


\begin{verbatim}

ClimateProcessingEngine

        |
        |
Configuration Layer
        |
        |
Input / Data Pipeline
        |
        |
Numerical Processing Engine
        |
        |
Statistical Distribution Engine
        |
        |
Result Pipeline
        |
        |
Monitoring System

\end{verbatim}



\subsection{Core Components}



\subsubsection{Configuration Layer}


Responsible for centralized management of:


\begin{itemize}

\item File paths
\item Climate variables
\item Processing parameters
\item Validation rules
\item Logging configuration

\end{itemize}



\subsubsection{Input Pipeline}


The input layer handles:


\begin{itemize}

\item Reading climate archives
\item Loading station blocks
\item Data assembly
\item Input validation

\end{itemize}



\subsubsection{Numerical Engine}


The numerical core performs:


\begin{itemize}

\item Window extraction
\item Statistical calculation
\item Distribution fitting
\item Parameter estimation
\item Model comparison

\end{itemize}



\subsubsection{Result Pipeline}


Responsible for:


\begin{itemize}

\item Output validation
\item Result generation
\item Zarr dataset writing

\end{itemize}



\subsubsection{Monitoring System}


Provides:


\begin{itemize}

\item Logging
\item Benchmarking
\item Checkpoint recovery
\item Performance monitoring

\end{itemize}



\section{Processing Workflow}


The complete workflow is:


\begin{enumerate}

\item Load raw climate datasets

\item Assemble station-based data blocks

\item Validate input data

\item Extract temporal windows

\item Fit statistical distributions

\item Select the best probability model

\item Validate results

\item Write output datasets

\item Store checkpoint information

\end{enumerate}




\section{Statistical Distribution Fitting}


ClimateProcessingEngine fits multiple probability distributions and selects
the optimal model using information criteria.



\subsection{Normal Distribution}


The probability density function is:


\begin{equation}

f(x)=
\frac{1}{\sqrt{2\pi\sigma^2}}
\exp
\left(
-\frac{(x-\mu)^2}{2\sigma^2}
\right)

\end{equation}



Parameters:


\begin{itemize}

\item $\mu$: Mean value

\item $\sigma$: Standard deviation

\end{itemize}



\subsection{Skew-Normal Distribution}



The skew-normal density is:


\begin{equation}

f(x)=
\frac{2}{\omega}
\phi
\left(
\frac{x-\xi}{\omega}
\right)
\Phi
\left(
\alpha
\frac{x-\xi}{\omega}
\right)

\end{equation}



Parameters:


\begin{itemize}

\item $\xi$: Location parameter

\item $\omega$: Scale parameter

\item $\alpha$: Shape parameter

\end{itemize}




\subsection{Bimodal Normal Distribution}



The bimodal model is defined as a mixture distribution:


\begin{equation}

f(x)=
w_1
N(x;\mu_1,\sigma_1^2)
+
w_2
N(x;\mu_2,\sigma_2^2)

\end{equation}



where:


\begin{itemize}

\item $w_1,w_2$: Mixture weights

\item $\mu_1,\mu_2$: Distribution means

\item $\sigma_1,\sigma_2$: Standard deviations

\end{itemize}



\subsection{Pearson Type III Distribution}



The Pearson Type III probability density function:


\begin{equation}

f(x)=
\frac{
1
}{
\Gamma(\alpha)\beta^\alpha
}
(x-\gamma)^{\alpha-1}
\exp
\left(
-\frac{x-\gamma}{\beta}
\right)

\end{equation}



Parameters:


\begin{itemize}

\item $\alpha$: Shape parameter

\item $\beta$: Scale parameter

\item $\gamma$: Location parameter

\end{itemize}




\section{Model Selection}


The engine evaluates candidate distributions using:



\subsection{AIC}


\begin{equation}

AIC=2k-2\ln(\hat{L})

\end{equation}



\subsection{Corrected AIC}


\begin{equation}

AIC_c
=
AIC+
\frac{2k(k+1)}
{n-k-1}

\end{equation}



\subsection{Bayesian Information Criterion}


\begin{equation}

BIC=
k\ln(n)-2\ln(\hat{L})

\end{equation}



where:


\begin{itemize}

\item $k$: Number of parameters

\item $n$: Number of observations

\item $\hat{L}$: Maximum likelihood estimate

\end{itemize}



\section{Bimodality Metrics}



\subsection{Ashman's D}


The separation between two Gaussian modes:


\begin{equation}

D=
\frac{
|\mu_2-\mu_1|
}
{
\sqrt{
\frac{\sigma_1^2+\sigma_2^2}{2}
}
}

\end{equation}



\subsection{Overlap Coefficient}



\begin{equation}

OVL=
\int_{-\infty}^{\infty}
\min
\{f_1(x),f_2(x)\}
dx

\end{equation}
\section{Installation}


\subsection{System Requirements}


Recommended environment:


\begin{itemize}

\item Python 3.10 or higher

\item Minimum RAM: 16 GB

\item Recommended RAM: 32 GB

\item Storage: 100 GB free space

\item CPU: 4 cores minimum

\item Recommended CPU: 8+ cores

\end{itemize}



\subsection{Clone Repository}


\begin{verbatim}

git clone https://github.com/AminFazlKazemi/ClimateProcessingEngine.git

cd ClimateProcessingEngine

\end{verbatim}



\subsection{Create Virtual Environment}



Windows:


\begin{verbatim}

python -m venv venv

venv\Scripts\activate

\end{verbatim}



Linux/macOS:


\begin{verbatim}

python3 -m venv venv

source venv/bin/activate

\end{verbatim}



\subsection{Install Dependencies}



\begin{verbatim}

pip install -r requirements.txt

\end{verbatim}




\section{Configuration}



The main configuration file is:


\begin{verbatim}

climatology_engine/config.yaml

\end{verbatim}



Example configuration:


\begin{verbatim}

paths:

  calendar_file:
      K:/Temp/calendar.txt

  checkpoint_file:
      checkpoint.csv

  output_dir:
      I:/climatology_output

  output_zarr_name:
      climatology_stationwise_final.zarr


years:

  start: 1369

  end: 1399


days: 366


window_days: 2



processing:

  block_size: 1000

  cores: 6

  use_parallel: false



variables:

  - tmin

  - tmean

  - tmax



logging:

  level: INFO

  log_file:
      climatology.log


\end{verbatim}



\section{Important Constants}



\begin{table}[h]

\centering

\begin{tabular}{lll}

\toprule

Constant & Value & Description \\

\midrule

WINDOW\_SIZE & 5 &
Two days before and after target day \\


MAX\_VALUES\_PER\_FIT & 155 &
Maximum fitting samples \\


MIN\_VALID\_VALUES & 5 &
Minimum valid observations \\


N\_OUTPUTS & 87 &
Number of output variables \\


VALID\_BEST\_DIST &
\{-1,0,1,2,3\} &
Distribution codes \\


\bottomrule

\end{tabular}

\end{table}



\section{Usage Guide}



\subsection{Running the Engine}



\begin{verbatim}

cd climatology_engine

python main.py

\end{verbatim}



The engine automatically detects existing checkpoints and resumes
unfinished processing.



\subsection{Starting a New Processing Run}



\begin{verbatim}

rm checkpoint.csv

python main.py

\end{verbatim}



\subsection{Analyzing Results}



After completion:


\begin{verbatim}

python analyze_distributions1.py

\end{verbatim}



\section{Python Example}



Load generated Zarr output:


\begin{verbatim}

import xarray as xr


ds = xr.open_zarr(
"climatology_stationwise_final.zarr"
)


print(list(ds.data_vars))


best_distribution = (
ds.tmin_best_dist
.isel(point=0)
.values
)


print(best_distribution.shape)

\end{verbatim}




\section{Output Data}



The engine produces:


\begin{itemize}


\item Best fitted distribution type

\item Distribution parameters

\item Statistical scores

\item Quality indicators

\item Station-wise climatological summaries

\item Zarr optimized datasets


\end{itemize}
\section{Performance Optimization}



ClimateProcessingEngine is designed for efficient processing of
large-scale climate datasets.



\subsection{Block-Based Processing}



Large datasets are divided into smaller blocks:


\begin{itemize}

\item Reduces memory consumption

\item Improves cache efficiency

\item Enables checkpoint recovery

\item Supports parallel execution

\end{itemize}



\subsection{Parallel Processing}



The engine supports:


\begin{itemize}

\item Multiprocessing

\item Concurrent execution

\item CPU-based parallel computation

\end{itemize}



The number of workers can be configured:


\begin{verbatim}

cores: 6

use_parallel: true

\end{verbatim}




\subsection{Memory Management}



Optimization strategies:


\begin{itemize}

\item Block-wise data loading

\item Garbage collection

\item Efficient NumPy arrays

\item Float32 memory optimization

\item Chunked Zarr storage

\end{itemize}




\subsection{Performance Recommendations}



Recommended practices:


\begin{enumerate}

\item Use SSD storage

\item Increase CPU cores when available

\item Adjust block size according to RAM

\item Disable unnecessary validation during production runs

\item Close memory-intensive applications

\end{enumerate}





\section{Project Structure}



The project organization:



\begin{verbatim}

ClimateProcessingEngine/

|

|-- climatology_engine/

|   |

|   |-- main.py

|   |-- constants.py

|   |-- config.yaml

|   |-- zarr_schema.py

|   |-- bimodal_normal.py

|   |-- calendar_tables.py

|   |-- runtime_tables.py

|

|   |-- io_pipeline/

|   |     |

|   |     |-- read_month_files.py

|   |     |-- assemble_block.py

|   |     |-- validate_block.py

|

|   |-- numerical_engine/

|   |     |

|   |     |-- window_engine.py

|   |     |-- distributions.py

|   |     |-- analyze_station.py

|   |     |-- merge_results.py

|

|   |-- orchestrator/

|   |     |

|   |     |-- process_block.py

|

|   |-- result_pipeline/

|   |     |

|   |     |-- validate_result.py

|   |     |-- write_block.py

|

|   |-- monitoring/

|         |

|         |-- benchmark.py

|         |-- checkpoint.py

|         |-- logger.py


|

|-- README.md

|-- README.tex

|-- LICENSE

|-- requirements.txt

|-- CITATION.cff


\end{verbatim}





\section{Dependencies}



Required Python packages:



\begin{verbatim}

numpy>=1.24.0

pandas>=2.0.0

xarray>=2023.0.0

dask>=2023.0.0

rasterio>=1.3.0

geopandas>=0.14.0

cartopy>=0.22.0

matplotlib>=3.7.0

scipy>=1.10.0

scikit-learn>=1.3.0

shapely>=2.0.0

netcdf4>=1.6.0

zarr>=2.15.0

numba>=0.57.0

pyyaml>=6.0

psutil>=5.9.0


\end{verbatim}




\section{Scientific Computing Stack}



The framework integrates:



\begin{table}[h]

\centering

\begin{tabular}{ll}

\toprule

Category & Libraries \\

\midrule

Numerical Computing &
NumPy, SciPy \\


Data Processing &
Pandas, Xarray, Dask \\


Geospatial &
Rasterio, GeoPandas, Cartopy, Shapely \\


Machine Learning &
Scikit-learn \\


Performance &
Numba \\


Visualization &
Matplotlib \\


Storage &
Zarr, NetCDF \\


\bottomrule

\end{tabular}

\end{table}
\section{Troubleshooting}



\subsection{File Not Found Error}



Common error:



\begin{verbatim}

FileNotFoundError:
No such file or directory

\end{verbatim}



Possible solutions:


\begin{itemize}

\item Check paths in config.yaml

\item Use absolute paths

\item Verify input files exist

\item Check read permissions

\end{itemize}




\subsection{Memory Error}



Example:


\begin{verbatim}

MemoryError:
Unable to allocate array

\end{verbatim}



Solutions:


\begin{itemize}

\item Reduce block size

\item Disable parallel processing

\item Reduce validation frequency

\item Increase available RAM

\end{itemize}





\subsection{Key Error}



Example:


\begin{verbatim}

KeyError:
variable_name

\end{verbatim}



Solutions:


\begin{itemize}

\item Check variable names

\item Verify input dataset schema

\item Update configuration variables

\end{itemize}





\section{Logging and Debugging}



The engine provides detailed logging:



\begin{itemize}

\item Console progress information

\item Log file records

\item Processing checkpoints

\item Error tracking

\end{itemize}



Available log levels:



\begin{verbatim}

DEBUG

INFO

WARNING

ERROR

\end{verbatim}





\subsection{Debugging Recommendations}



\begin{enumerate}

\item Enable DEBUG logging

\item Check climatology.log

\item Run a small test block

\item Verify input data quality

\item Use Python debugger if necessary

\end{enumerate}





\section{Contributing}



Contributions are welcome.



Contribution workflow:



\begin{enumerate}


\item Fork the repository


\item Create a new branch:


\begin{verbatim}

git checkout -b feature-name

\end{verbatim}



\item Commit changes:


\begin{verbatim}

git commit -m "Add new feature"

\end{verbatim}



\item Push changes:


\begin{verbatim}

git push origin feature-name

\end{verbatim}



\item Submit a Pull Request


\end{enumerate}





\section{Citation}



If you use ClimateProcessingEngine in scientific research,
please cite this project:



\begin{quote}

Amin Fazl Kazemi.

\textbf{ClimateProcessingEngine:
A Python Framework for Large-Scale Climate Data Processing,
Statistical Distribution Fitting and Scientific Analysis.}

GitHub Repository:

\url{https://github.com/AminFazlKazemi/ClimateProcessingEngine}

\end{quote}




\section{License}



This project is released under the MIT License.



See:


\begin{verbatim}

LICENSE

\end{verbatim}



for complete license information.





\section{Contact}



\textbf{Amin Fazl Kazemi}



GitHub:


\url{https://github.com/AminFazlKazemi}



LinkedIn:


\url{https://linkedin.com/in/aminfazlkazemi}



Email:


\url{aminfazlkazemi@gmail.com}





\section*{Acknowledgements}



ClimateProcessingEngine is built upon the open-source scientific Python ecosystem
and benefits from contributions by the global climate and data science community.





\section*{Final Note}



If this project is useful for your climate research,
please consider giving it a star on GitHub.



\begin{center}

\textbf{ClimateProcessingEngine}

\end{center}



\end{document}